% GENERATED FILE. Edit canonical lessons, then run npm run generate:books.
% canonical-source-hash: fnv1a64:4260173fa491e315
% canonical-lessons: ML-C424-pariskarikkuka, ML-C424-niyamanam, ML-C424-janadhipatyam, ML-C424-adhikaram, ML-C424-netavu

\chapter{Posts and Reforms}
\label{ch:ml-posts-and-reforms}
\hlchaptermodality{424}

\begin{chapteropening}
\textbf{By the end of this chapter:} \emph{I can say an appointment, democracy, authority, a leader and to reform.}

Complete the last lesson of chapter 424: I can say an appointment, democracy, authority, a leader and to reform.
\end{chapteropening}

\section[\texorpdfstring{\ml{പരിഷ്കരിക്കുക} (pariṣkarikkuka)}{പരിഷ്കരിക്കുക (pariṣkarikkuka)}]{\ml{പരിഷ്കരിക്കുക} (pariṣkarikkuka) — to reform; to revise}

\label{lesson:ML-C424-pariskarikkuka}



\begin{quote}
Before the new one: say the Malayalam for to amend, then the Malayalam for an oath.
\end{quote}

\subsection*{You'll want to know: \ml{പരിഷ്കരിക്കുക}}
\textbf{\ml{പരിഷ്കരിക്കുക}} — \emph{pariṣkarikkuka} — \textquotedblleft{}to reform; to revise\textquotedblright{}.

\begin{quote}
\ml{സർക്കാർ} \ml{നികുതി} \ml{നിയമം} \ml{പരിഷ്കരിച്ചു}.
\end{quote}

\begin{quote}
\emph{sarkkār nikuti niyamaṁ pariṣkariccu.}
\end{quote}

\begin{quote}
The government reformed the tax law.
\end{quote}

Past: \ml{പരിഷ്കരിച്ചു}, \emph{pariṣkariccu}.

The noun is \ml{പരിഷ്കാരം}, a reform.

\begin{cousinweb}
\textbf{\ml{പരിഷ്കരിക്കുക}} (\emph{pariṣkarikkuka}) is built from Sanskrit \emph{pariṣkāra}, \textquotedblleft{}an ornament\textquotedblright{} (CDIAL 7894), with the Malayalam verb ending -\ml{ിക്കുക}. To reform a thing is to dress it anew.
\end{cousinweb}

\begin{grammarlens}[title={What you've built}]
One more word for giving reasons and opinions.
\end{grammarlens}

\subsection*{Your turn}
\begin{itemize}
\raggedright
  \item \emph{Say it:} \emph{pariṣkarikkuka}
  \item \emph{Say it:} \emph{pariṣkarikkuka}, once more
  \item \emph{Say it:} \emph{sarkkār nikuti niyamaṁ pariṣkariccu.} — The government reformed the tax law.
  \item \emph{Recall:} say the Malayalam for to amend, then the Malayalam for an oath, then say \emph{pariṣkarikkuka} again
\end{itemize}

\begin{tcolorbox}[breakable,colback=teal!4,colframe=teal!35!black,title={Before you move on}]
What does \ml{പരിഷ്കരിക്കുക} mean? (\textquotedblleft{}To reform; to revise\textquotedblright{}.) Say it once more.
\end{tcolorbox}
\section[\texorpdfstring{\ml{നിയമനം} (niyamanaṁ)}{നിയമനം (niyamanaṁ)}]{\ml{നിയമനം} (niyamanaṁ) — an appointment (to a post)}

\label{lesson:ML-C424-niyamanam}



\begin{quote}
Before the new one: say the Malayalam for an oath, then the Malayalam for to reform.
\end{quote}

\subsection*{You'll want to know: \ml{നിയമനം}}
\textbf{\ml{നിയമനം}} — \emph{niyamanaṁ} — \textquotedblleft{}an appointment (to a post)\textquotedblright{}.

\begin{quote}
\ml{പുതിയ} \ml{അധ്യാപകരുടെ} \ml{നിയമനം} \ml{അടുത്ത} \ml{മാസം} \ml{നടക്കും}.
\end{quote}

\begin{quote}
\emph{putiya adhyāpakaruṭe niyamanaṁ aṭutta māsaṁ naṭakkuṁ.}
\end{quote}

\begin{quote}
The appointment of the new teachers will take place next month.
\end{quote}

A post filled by an official order; a meeting time is a different word.

\begin{cousinweb}
\textbf{\ml{നിയമനം}} (\emph{niyamanaṁ}) is borrowed from Sanskrit \emph{niyamana}, \textquotedblleft{}subduing, restraining\textquotedblright{} (Monier-Williams), from the same root as \ml{നിയമം} (CDIAL 7250); Malayalam uses it for fixing a person in a post.
\end{cousinweb}

\begin{grammarlens}[title={What you've built}]
One more word for giving reasons and opinions.
\end{grammarlens}

\subsection*{Your turn}
\begin{itemize}
\raggedright
  \item \emph{Say it:} \emph{niyamanaṁ}
  \item \emph{Say it:} \emph{niyamanaṁ}, once more
  \item \emph{Say it:} \emph{putiya adhyāpakaruṭe niyamanaṁ aṭutta māsaṁ naṭakkuṁ.} — The appointment of the new teachers will take place next month.
  \item \emph{Recall:} say the Malayalam for an oath, then the Malayalam for to reform, then say \emph{niyamanaṁ} again
\end{itemize}

\begin{tcolorbox}[breakable,colback=teal!4,colframe=teal!35!black,title={Before you move on}]
What does \ml{നിയമനം} mean? (\textquotedblleft{}An appointment (to a post)\textquotedblright{}.) Say it once more.
\end{tcolorbox}
\section[\texorpdfstring{\ml{ജനാധിപത്യം} (janādhipatyaṁ)}{ജനാധിപത്യം (janādhipatyaṁ)}]{\ml{ജനാധിപത്യം} (janādhipatyaṁ) — democracy}

\label{lesson:ML-C424-janadhipatyam}



\begin{quote}
Before the new one: say the Malayalam for to reform, then the Malayalam for an appointment.
\end{quote}

\subsection*{You'll want to know: \ml{ജനാധിപത്യം}}
\textbf{\ml{ജനാധിപത്യം}} — \emph{janādhipatyaṁ} — \textquotedblleft{}democracy\textquotedblright{}.

\begin{quote}
\ml{ജനാധിപത്യത്തിൽ} \ml{ജനങ്ങളാണ്} \ml{സർക്കാരിനെ} \ml{തിരഞ്ഞെടുക്കുന്നത്}.
\end{quote}

\begin{quote}
\emph{janādhipatyattil janaṅṅaḷāṇŭ sarkkārine tiraññeṭukkunnatŭ.}
\end{quote}

\begin{quote}
In a democracy it is the people who choose the government.
\end{quote}

The adjective is \ml{ജനാധിപത്യ}: \ml{ജനാധിപത്യ} \ml{രാജ്യം}, a democratic country.

\begin{cousinweb}
\textbf{\ml{ജനാധിപത്യം}} (\emph{janādhipatyaṁ}) is a modern word built from Sanskrit \emph{jana}, \textquotedblleft{}people\textquotedblright{} (CDIAL 5098), and \emph{ādhipatya}, \textquotedblleft{}lordship, sovereignty\textquotedblright{} (Monier-Williams): rule by the people.
\end{cousinweb}

\begin{grammarlens}[title={What you've built}]
One more word for giving reasons and opinions.
\end{grammarlens}

\subsection*{Your turn}
\begin{itemize}
\raggedright
  \item \emph{Say it:} \emph{janādhipatyaṁ}
  \item \emph{Say it:} \emph{janādhipatyaṁ}, once more
  \item \emph{Say it:} \emph{janādhipatyattil janaṅṅaḷāṇŭ sarkkārine tiraññeṭukkunnatŭ.} — In a democracy it is the people who choose the government.
  \item \emph{Recall:} say the Malayalam for to reform, then the Malayalam for an appointment, then say \emph{janādhipatyaṁ} again
\end{itemize}

\begin{tcolorbox}[breakable,colback=teal!4,colframe=teal!35!black,title={Before you move on}]
What does \ml{ജനാധിപത്യം} mean? (\textquotedblleft{}Democracy\textquotedblright{}.) Say it once more.
\end{tcolorbox}
\section[\texorpdfstring{\ml{അധികാരം} (adhikāraṁ)}{അധികാരം (adhikāraṁ)}]{\ml{അധികാരം} (adhikāraṁ) — authority; power}

\label{lesson:ML-C424-adhikaram}



\begin{quote}
Before the new one: say the Malayalam for an appointment, then the Malayalam for democracy.
\end{quote}

\subsection*{You'll want to know: \ml{അധികാരം}}
\textbf{\ml{അധികാരം}} — \emph{adhikāraṁ} — \textquotedblleft{}authority; power\textquotedblright{}.

\begin{quote}
\ml{അവന്} \ml{അത്} \ml{ചെയ്യാൻ} \ml{അധികാരം} \ml{ഇല്ല}.
\end{quote}

\begin{quote}
\emph{avanŭ atŭ ceyyān adhikāraṁ illa.}
\end{quote}

\begin{quote}
He has no authority to do that.
\end{quote}

An official is \ml{അധികാരി}, one with authority.

\begin{cousinweb}
\textbf{\ml{അധികാരം}} (\emph{adhikāraṁ}) is borrowed from Sanskrit \emph{adhikāra}, \textquotedblleft{}authority; government, rule\textquotedblright{} (Monier-Williams).
\end{cousinweb}

\begin{grammarlens}[title={What you've built}]
One more word for giving reasons and opinions.
\end{grammarlens}

\subsection*{Your turn}
\begin{itemize}
\raggedright
  \item \emph{Say it:} \emph{adhikāraṁ}
  \item \emph{Say it:} \emph{adhikāraṁ}, once more
  \item \emph{Say it:} \emph{avanŭ atŭ ceyyān adhikāraṁ illa.} — He has no authority to do that.
  \item \emph{Recall:} say the Malayalam for an appointment, then the Malayalam for democracy, then say \emph{adhikāraṁ} again
\end{itemize}

\begin{tcolorbox}[breakable,colback=teal!4,colframe=teal!35!black,title={Before you move on}]
What does \ml{അധികാരം} mean? (\textquotedblleft{}Authority; power\textquotedblright{}.) Say it once more.
\end{tcolorbox}
\section[\texorpdfstring{\ml{നേതാവ്} (nētāvŭ)}{നേതാവ് (nētāvŭ)}]{\ml{നേതാവ്} (nētāvŭ) — a leader}

\label{lesson:ML-C424-netavu}



\begin{quote}
Before the new one: say the Malayalam for democracy, then the Malayalam for authority.
\end{quote}

\subsection*{You'll want to know: \ml{നേതാവ്}}
\textbf{\ml{നേതാവ്}} — \emph{nētāvŭ} — \textquotedblleft{}a leader\textquotedblright{}.

\begin{quote}
\ml{ഗ്രാമത്തിലെ} \ml{ആളുകൾ} \ml{അവനെ} \ml{നേതാവായി} \ml{തിരഞ്ഞെടുത്തു}.
\end{quote}

\begin{quote}
\emph{grāmattile āḷukaḷ avane nētāvāyi tiraññeṭuttu.}
\end{quote}

\begin{quote}
The people of the village chose him as their leader.
\end{quote}

Plural: \ml{നേതാക്കൾ}, \emph{nētākkaḷ}.

Plural: \ml{നേതാക്കൾ}.

\begin{cousinweb}
\textbf{\ml{നേതാവ്}} (\emph{nētāvŭ}) is borrowed from Sanskrit \emph{netṛ}, \textquotedblleft{}leading, a leader\textquotedblright{} (Monier-Williams), whose nominative is \emph{netā}, with the Malayalam ending -\ml{വ്}.
\end{cousinweb}

\begin{grammarlens}[title={What you've built}]
One more word for giving reasons and opinions.
\end{grammarlens}

\subsection*{Your turn}
\begin{itemize}
\raggedright
  \item \emph{Say it:} \emph{nētāvŭ}
  \item \emph{Say it:} \emph{nētāvŭ}, once more
  \item \emph{Say it:} \emph{grāmattile āḷukaḷ avane nētāvāyi tiraññeṭuttu.} — The people of the village chose him as their leader.
  \item \emph{Recall:} say the Malayalam for democracy, then the Malayalam for authority, then say \emph{nētāvŭ} again
\end{itemize}

\begin{tcolorbox}[breakable,colback=teal!4,colframe=teal!35!black,title={Before you move on}]
What does \ml{നേതാവ്} mean? (\textquotedblleft{}A leader\textquotedblright{}.) Say it once more.
\end{tcolorbox}
